Análisis
La distribución de los saldos por rango de edad durante 2025 muestra que la mayor concentración de recursos se encuentra en los clientes con edades comprendidas entre
20 y 69 años
. En estos grupos, el saldo agregado oscila entre aproximadamente
Q2.15 millones
y
Q2.70 millones
, evidenciando una distribución relativamente homogénea entre las diferentes etapas de la vida laboral.
Por el contrario, los grupos ubicados en los extremos de la distribución, correspondientes a clientes de
15 a 19 años
y de
70 a 74 años
, registran los menores niveles de saldo, con montos aproximados entre
Q464 mil
y
Q790 mil
. La diferencia respecto a los grupos de mayor concentración resulta considerable, lo que sugiere una menor participación de estos rangos de edad en el volumen total de recursos administrados por el banco.
Este comportamiento es consistente con el ciclo de vida financiero de los clientes, ya que las edades comprendidas entre los 20 y 69 años suelen concentrar una mayor participación en el mercado laboral y, por consiguiente, una mayor capacidad para generar y mantener saldos en sus cuentas. No obstante, esta interpretación constituye una hipótesis de negocio y requeriría información adicional, como el número de clientes por rango de edad, el saldo promedio por cliente y otras variables socioeconómicas, para ser validada.